\section{Discussion}
\label{sec:discussion}

\subsection{Interpretation}

The coverage result is structural, following from what each evidence class can
physically observe, and it was obtained
with calibration held fixed across all scenarios. The implication for the
literature is narrow but practical. An accuracy reported against
localisation-visible attacks is an accuracy over that subset of the attack
space. EMBN-MANET's threat model is confined to false-location behaviour by
design, so its figure is correct within its scope and we do not dispute it.
What we would ask of any such report is one additional line: which attack
classes the evidence can observe at all.

The isolation result is the paper's strongest: confirmation-gated isolation
withdrew 3.9--5.0\% of honest-beacon service against 74.8\% for the strongest
immediate policy, and the ordering held under attack, under shadowing and when
enforced. Its cost is the time an attacker remains active before
confirmation, a median of 15--27.5\,s, or 125--150\,s for selective forwarding
under shadowing. Whether that trade is favourable depends on the prior
probability of attack and the cost of a lost beacon: gating suits rare attacks
and sparse beacon coverage, because $\Phi$ accumulates continuously while
exposure is bounded by the confirmation delay, and immediate isolation suits
frequent attacks and high redundancy (Table~\ref{tab:guide}). The bound on
exposure holds only if release cannot be triggered by missing evidence, and in
our data the safe release rule made isolation on confirmation close to
permanent. We suggest that $\Phi$ or a comparable measure be reported
routinely, since a false-positive rate says nothing about how long errors
persist.

\begin{table}[!t]
\caption{Isolation policy by operating condition, as suggested by the results.}
\label{tab:guide}
\centering
\footnotesize
\begin{tabular}{@{}>{\raggedright\arraybackslash}p{4.4cm}>{\raggedright\arraybackslash}p{2.2cm}>{\raggedright\arraybackslash}p{5.4cm}@{}}
\toprule
Condition & Policy & Basis in the results \\
\midrule
Attacks rare, few redundant beacons & E$'$ & About 4--5\% of honest-beacon service withdrawn; attacker active for the confirmation delay \\
Attacks frequent, many redundant beacons & K$_1$ (D$_1$) or K$_2$ & Removes attackers soonest; K$_2$ withdraws 22--24\% of honest-beacon service against 60--75\% for D$_1$ \\
False isolation must be near zero & K$_4$ or K$_5$ & 0--0.8\% of honest-beacon service withdrawn, for 34--53\,s of attacker exposure \\
Forwarding-only attackers expected & E$'$, not E & E re-admits forwarding attackers once isolation silences their evidence \\
Sparse network & E$'$ with peer-relative behavioural scoring & No false isolation observed, but forwarding attackers confirmed in 4 of 10 seeds \\
Heavy shadowing & E$'$ & False isolation stays low; expect 2--3\,min to confirm forwarding attacks \\
\bottomrule
\end{tabular}
\end{table}

Two practical limits apply to this guidance. First, $\Phi$ counts lost
beacon service time, not metres of localisation error. Our runs did not
recompute the position estimates of nodes near a wrongly isolated beacon, so
the paper cannot convert one into the other; doing so needs the node density
around each beacon and the localisation algorithm in use. Second, the
pipeline falsely confirmed about one honest beacon in ten in the dense
network, and under E$'$ such isolations were almost never reversed. A
long-running deployment would lose beacons slowly, so it needs a planned
path back, such as the probationary release of Section~\ref{sec:future}.

The release finding is not specific to our pipeline. Any scheme that
restores trust as bad evidence fades has the same exposure, because an
isolated node produces no evidence at all. CONFIDANT's re-integration of
nodes whose reputation recovers over time~\cite{buchegger2002confidant} and
OCEAN's faulty-list timeout~\cite{bansal2003ocean} are examples: a forwarding
attacker excluded from routes stops generating the observations that would
keep its reputation low. Earlier critiques of redemption assumed that the
attacker waits out a timer or behaves well to rebuild its
score~\cite{sen2010survey,sun2006attacks}; here re-admission needs no action
by the attacker at all, which is why a defence against on--off behaviour
does not by itself prevent it.

\subsection{Limits of the evidence}

Whether fusion helped was decided by the attack family, not by the
combination rule. The same detector diluted a disjunctive rule when the
branches' sensitivity ratio was 3.63 and matched or slightly exceeded it at
1.87, as the $1+\sqrt2$ bound predicts. The baseline comparison is the same
result at the level of whole detectors: in the small dense network two
calibrated single-evidence detectors run in parallel matched the pipeline,
while in the sparse and 50-node networks they withdrew 11--24\% of
honest-beacon service against 0--3\%, through peer-relative scoring in the
first case and fusion in the second. When fusion helps, it does so by
suppressing single-branch false alarms, at the price of sensitivity near the
threshold. Branch sensitivity to an attack family can be measured before
fusion is applied, and no system we surveyed reports checking it; because an
adaptive adversary could concentrate its misbehaviour on one axis, a
single-branch path to confirmation should stay alongside the fused statistic.

The aggregation stage changed no verdict in the ablation (Online Resource~1,
Sect.~S6), so a system using only direct per-observer trust would reproduce
every reported finding. It remains as an unevaluated extension for conditions
our campaign barely samples: observers who disagree, noisy direct trust and
recommenders who lie. Neither isolation policy produced a defensible
throughput benefit, possibly because AODV finds alternative routes readily in
a dense 25-node topology~\cite{sirajuddin2021tbsmr,kavitha2019india,buksh2024trust};
on our evidence the claim that isolation restores network performance is
unsupported. Throughput says nothing about A1, which has no packet-level
footprint; localisation error is the appropriate metric for that attack.

\section{Threats to Validity}
\label{sec:threats}

\textbf{Construct validity.} The pipeline runs over logged simulator traces:
recommendation messages were not simulated (Section~\ref{sec:cost}), only the
resulting isolation schedules were enforced, and behavioural monitoring is
node-global rather than per-observer, which may explain the null aggregation
ablation. $\Phi$ for the 30-seed campaign assumes that an isolation does not
change later verdicts. This is exact for the first isolation of a
never-releasing policy but not for E's release times
(Section~\ref{sec:release}). Enforced replay on ten clean seeds moved neither
policy's $\Phi$ distinguishably from zero, and three closed-loop rounds left
E$'$ at a fixed point; iteration approximates an online detector but is not
one. E$'$ was designed after the enforced runs were examined, so only its
sparse and 50-node results (no attacker released, clean false isolation of 0
and 2.6\%) are out-of-sample. In the sparse calibration runs every usable
honest behavioural window forwarded all its transit packets, so the
behavioural scale collapsed to its floor; this makes the branch maximally
sensitive and cannot explain the low sparse detection rate, but the sparse
behavioural figures should be read with it in mind.

\textbf{Internal validity.} Thresholds were calibrated on held-out clean runs
and fixed, which avoids tuning to the test scenarios but gives no scenario its
optimal operating point. The recovery baseline window is short relative to
throughput variance, which is why throughput debt is excluded.

\textbf{External validity.} Detection and the false-isolation result survive
4--8\,dB shadowing, at a cost in latency for forwarding attacks, but the
channel has no multipath or obstruction model. Attackers do not collude or
adapt, mobility is synthetic, and only the hybrid attack's intensity is swept.
The throughput nulls rest on the small dense field alone. Memory limited the
network size: a 50-node run needs about 3\,GB, and a 100-node configuration
did not fit. Three monitored beacons per run, six at 50 nodes, make
rate-based metrics coarse. We do not claim generality beyond this
configuration.

\textbf{Comparative validity.} The baselines re-implement the evidence class
and decision rule of Liu--Ning--Du alert counting and watchdog/CONFIDANT
reputation on our counters and residual, with thresholds set by rules we
chose; a different defensible rule would move their operating points. Policies
D, D$_w$ and D$_1$ are our own constructions. In Section~\ref{sec:frontier}
the persistence settings share D$_1$'s replay assumption, and baseline
exposure in the sparse and 50-node networks is estimated from detection counts
and latencies. The comparison with EMBN-MANET is analytical.

\textbf{Statistical conclusion validity.} Coverage, latency, false positives
and $\Phi$ rest on 30 seeds with paired bootstrap intervals; throughput,
recovery, localisation error and the ablation on five; the sparse and 50-node
campaigns and the sweeps on ten, so the one-seed fusion gain is reported as
consistent with the condition, not as a benefit. About forty paired intervals
are reported without multiplicity correction. The central comparisons
(immediate policies against E or E$'$) lie hundreds of node-seconds from zero;
intervals that only just exclude or touch zero, such as several 50-node
baseline comparisons, are indicative.

\section{Conclusion}
\label{sec:conclusion}

Across 780 ns-3 runs, single-evidence detection of malicious beacons was blind
to attacks outside its evidence class, and collecting both behavioural and
localisation evidence removed the observed blind spots. Confirmation-gated
isolation withdrew 3.9\% of honest-beacon service on a clean network against
74.8\% for the strongest immediate policy, at a median cost of 15--27.5\,s of
attacker activity, and the confirmation persistence alone traced the
trade-off between the two. Releasing a beacon when its verdict clears
re-admitted forwarding attackers, whose behavioural evidence disappears while
they are isolated; requiring that evidence to return prevented every
re-admission and raised false isolation only to 5.0\%. These results held
under 4 and 8\,dB shadowing, with the price moving to latency. A condition on
the branches' relative sensitivity predicted when Stouffer fusion diluted or
matched a disjunctive rule, and explains why two calibrated literature
baselines run side by side matched the pipeline in a small dense network but
withdrew 11--24\% of honest-beacon service, against 0--3\%, in sparse and
larger ones. The aggregation stage changed no verdict, and no isolation policy
produced a measurable throughput gain. In a literature where reported
accuracies cluster between 90 and 99\%, the more informative questions are
which attacks a detector can see and how long its mistakes last.

\subsection{Future work}
\label{sec:future}

Three directions follow from these results. The first is to run the detector
inside the simulator, so that every isolation and release feeds back into
later verdicts. That would test the evidence-aware release rule without the
replay assumption, and it would allow a probationary release in which a
released beacon is isolated again on its first malicious window. The second
is to test the findings in harder settings: networks of 100 nodes or more,
trace-based mobility, colluding or adaptive attackers, and enforced
isolation in sparse topologies, where isolation is most likely to affect
delivery. Full implementations of EMBN-MANET and of a behavioural
trust-routing scheme belong in the same study. The third is to speed up
behavioural confirmation, which dominated latency under shadowing. A second
behavioural observable, or a confidence weight that adapts to the channel,
could shorten the time to confirm forwarding attacks.

\section*{Acknowledgements}

The authors thank the developers and maintainers of the ns-3 network
simulator~\cite{ns3}, whose simulator and AODV, IEEE 802.11 and propagation
models this study builds on, and the developers of the open-source Python
scientific stack used for the analysis.

\section*{Declarations}

\textbf{Funding.} No funding was received for this work.

\textbf{Competing interests.} The authors declare no competing interests.

\textbf{Author contributions.} A.V.\ designed and implemented the simulator,
verification pipeline and experiments, and drafted the manuscript. D.V.C.\
supervised the work, contributed to the experimental design and revised the
manuscript. Both authors read and approved the final version.

\textbf{Supplementary information.} Online Resource~1 contains the
intensity-sweep tables, supplementary figures and secondary analyses cited in
the text.

\textbf{Data and code availability.} Simulation sources, the analysis
pipeline, the scripts that run every campaign and regenerate every data
figure, and the result tables, per-window verdicts and isolation schedules
for every campaign are available at
\url{https://github.com/kraghu356/manet-beacon-trust}. The raw per-run
simulator traces are available from the corresponding author
on reasonable request. No third-party dataset was used: every data point in
this paper was generated by the authors' simulations.

\textbf{Use of AI tools.} An AI assistant (Claude, Anthropic) was used for
code analysis.

\textbf{Ethics approval.} Not applicable. This study involves neither human
participants nor animals.
